% Volume VI: Adversaries, Platforms, and Automated Judgment
% Part XII. Platforms and Automated Publicity
% Chapter 66: Moderation as Compressed Adjudication
% Spine: proof-spines/ch066-spine.md

\chapter{Moderation as Compressed Adjudication}
\label{ch:ch066}


At scale, moderation becomes \textbf{compressed adjudication}. Platforms cannot investigate each case as a court or full administrative process would; instead they substitute rules, classifiers, queues, and confidence thresholds for individualized inquiry. The governing problem is not whether error can be eliminated, but how judgment is allocated under constrained attention and asymmetric error costs.

This means that moderation architecture already decides which cases are seen first, which kinds of evidence count, which thresholds trigger action, and which mistakes are treated as tolerable. Compression is therefore substantive rather than merely logistical. The chapter's preface frames moderation as a system that formalizes public judgment under scarcity: procedural values are not absent, but translated into queue design, automated triage, and threshold-setting that decide at scale what full adjudication would decide in detail.
